\chapter{Théorie abstraite de la réécriture}
\label{chp.reecr}

\minitoc

\lettrine{L}{e} chapitre précédent utilisait par deux fois une relation
$\reecr{}$ pour réécriture les arbres de preuve. Dans la suite de cette
partie, on aura souvent besoin de relation de réécriture (qu'on
appellera aussi relation de réduction, suivant le contexte), et on se
propose donc d'en étudier la théorie générale en amont. Ce chapitre est
donc dédié à la théorie abstraite de la réécriture, portant sur les
système abstraits de réécriture (\foreignexpr{Abstract Rewriting Systems},
abrégé en ARS).

Le principal objectif du chapitre est de donner les outils nécessaires
pour étudier efficacement les systèmes de réécriture donnés dans les
prochaines chapitres. A FAIRE

On pourra se référer à \cite{TermRewrAllThat} pour une étude plus
approfondie des systèmes de réécriture.

\section{\'Etude des ARS}

Le cadre des ARS est particulièrement général~: on s'intéresse en soi à
n'importe quel ensemble muni d'une relation binaire. Dans les faits, c'est
le choix des propriétés d'intérêt ainsi que des objets concrets auxquels on
applique la théorie qui donne à ce cadre son sens. Ainsi, on peut déjà citer
le cas de $(\ProofLJ(\Sign),\reecr{})$ ou de $(\ProofLK(\Sign),\reecr{})$.
Abordons donc, dans un premier temps, les constructions élémentaires
fixant un peu plus le cadre de travail.

\subsection{Premières définitions et propriétés}

\noindent
Commençons par la définition la plus élémentaire, déjà évoquée avant.

\begin{definition}[Système de réécriture abstrait]
    On appelle système de réécriture abstrait (abrégé $\ARS$) une paire
    $(A,\varTo)$ où $\mathord{\varTo}\subseteq A^2$ est une relation
    binaire.
\end{definition}

L'exemple sur lequel il est bon de fixer son intuition pour appréhender les
ARS est celui d'un système de calcul (par exemple, on pourrait considérer
des configurations de machines de Turing avec la relation d'action par une
machine donnée). Dans un système de calcul, les
éléments sont des programme, c'est-à-dire des suites de symboles, et la
relation de réécriture permet de faire avancer l'exécution du programme.
Ainsi, les ARS ont beau être un cadre très large, on va tout de suite
introduire des propriétés précises qui sont d'intérêt ici, car nous voulons
décrire des comportements précis parmi ces ensembles munis d'une relation.
Les premières propriétés à introduire sont parfaitement classiques, et
mènent avant tout à pouvoir construire une relation d'équivalence.

\begin{definition}[Propriétés d'un ARS]
    Soit $(A,\varTo)$ un ARS. On dira que $\varTo$ est
    \begin{itemize}
    \item réflexive quand $\forall a \in A, a \varTo a$
    \item irréflexive quand $\forall a \in A, a\nvarTo a$
    \item symétrique quand
      $\forall a\,b\in A, a\varTo b \implies b \varTo a$
    \item transitive quand
      $\forall a\,b\,c\in A, a \varTo b \land b \varTo c \implies
      a \varTo c$
    \end{itemize}
\end{definition}

Ces propriétés sont importantes, mais contrairement au cas général, nous
voulons ici traiter de relations qui ne sont pas, par défaut, des relations
d'équivalence. En reprenant l'analogie d'un système de calcul, si pour deux
programmes $p,q$ la relation $p\varTo q$ signifie que $p$ devient $q$ lors de
son exécution, il nous est nécessaire de garder disponibles deux points de
vue concurrents~:
\begin{itemize}
\item le point de vue de l'exécution, dans lequel on a les deux programmes
  $p$ et $q$, distincts, qui sont reliés par le fait de s'exécuter l'un en
  l'autre~;
\item le point de vue sémantique, dans lequel le programme est avant tout une
  représentation de son résultat, ce qui correspond à l'identification
  $p = q$ entre deux programmes s'exécutant l'un en l'autre.
\end{itemize}

Pour pouvoir rendre compte de ses deux points de vue, on conserve toujours la
relation $\varTo$ mais on introduit aussi les meilleures approximations de
cette relation pour avoir une relation symétrique, transitive, réflexive\ldots

\begin{definition}[Clôture réflexive, symétrique, transitive]
    Soit $(A,\varTo)$ un ARS. On définit les relations suivantes~:
    \begin{itemize}
    \item $\varToR$ est la plus petite relation réflexive contenant $\varTo$~;
    \item $\varToS$ est la plus petite relation symétrique contenant $\varTo$~;
    \item $\varToT$ est la plus petite relation transitive contenant $\varTo$~;
    \item $\varToRt$ est la plus petite relation réflexive et transitive
      contenant $\varTo$~;
    \item $\varToRst$ est la plus petite relation d'équivalence contenant
      $\varTo$.
    \end{itemize}
    On appelle ces relations les clôtures (respectivement réflexive, symétrique,
    transitive, réflexive et transitive, et réflexive symétrique et transitive)
    de $\varTo$.
\end{definition}

\begin{remark}
    On peut se référer aux \cref{exo.tClot}, \cref{exo.rtClot}, \cref{exo.sClot}
    et \cref{exo.rstClot} pour une étude de ces clôtures de relations.
\end{remark}

\begin{example}
    Prenons comme système, pour une machine de Turing (potentiellement non
    déterministe) $\varTMM$ fixée, l'ensemble des configurations sur $\varTMM$.
    La relation $\varTo$ relie naturellement une configuration à une
    configuration qui la suit sous l'action de $\varTMM$. Remarquons que si
    la configuration comporte un état d'arrêt, alors elle est en relation avec
    aucune configuration. Analysons maintenant les différentes clôtures pour cette
    relation~:
    \begin{itemize}
        \item la clôture réflexive relie une configuration $C$ aux configurations
        $C'$ atteignables en au plus un temps de calcul~;
        \item la clôture symétrique relie une configuration $C$ aux configurations
        $C'$ qui diffèrent d'un temps de calcul (dans les deux sens)~;
        \item la clôture transitive relie une configuration $C$ à toutes les
        configurations $C'$ qu'elle peut atteindre dans le futur~;
        \item la clôture réflexive et transitive relie une configuration $C$ à
        toutes les configurations $C'$ qu'elle peut atteindre dans le futur au sens
        large~;
        \item la clôture d'équivalence relie une configuration $C$ aux configurations
        $C'$ atteignables depuis $C$ en avançant d'une étape de calcul, puis en
        \og remontant le temps\fg si nécessaire, et ainsi de suite un nombre fini de
        fois.
    \end{itemize}
\end{example}

\noindent
On donne aussi quelques opérations élémentaires sur les relations.

\begin{definition}[Transposée]
    Soit $(A,\varTo)$ un ARS. On appelle transposée de $\varTo$ la relation
    \[a \varToTr b \defeq b \varTo a\]
    On appelle composition de deux relations $\varRR,\varRR'$ sur $A$ la
    relation
    \[a\mathrel{(\mathord{\varRR}\circ\mathord{\varRR'})}c \defeq
    \exists b \in A, a \varRR b \land b \varRR c\]
    On définit l'itérée d'une relation $\varRR$ par
    \[\begin{cases}
        \varRR^0 \defeq \compr{(a,a)}{a \in A}\\
        \mathord{\varRR^{n+1}}\defeq \mathord{\varRR^n}\circ\mathord{\varRR}
    \end{cases}\]
    Enfin, on définit les plages d'itérées d'une relation $\varRR$ par
    \[\mathord{\varRR^{\leq n}}\defeq \bigcup_{k \leq n} \mathord{\varRR^k}\]
\end{definition}

Avec ces constructions élémentaires, on peut finalement commencer à exprimer des
propriétés réellement intéressantes évoquant l'exécution dans des systèmes de
calcul.

\subsection{Confluence}

Le premier point à aborder est celui de la cohérence des résultats d'un calcul.
Prenons par exemple le cas des configurations de machines de Turing. Rappelons
qu'une machine de Turing non déterministe est une machine dans laquelle la
table de transition est une relation~: ainsi, depuis une configuration $C$
donnée, on dispose d'un ensemble $\varTMM\opTM C$ de configurations suivantes.
Dans un tel cas, on pourrait imaginer qu'un calcul \og s'effectue correctement\fg
pour une machine de Turing lorsque, pour n'importe quelles transitions qu'on
emprunte depuis la configuration initiale $\varqZ\varwd\toInf$, on arrive à la même
configuration $\varqP{f}\varwv\toInf$. Ainsi, le non déterminisme de la machine de
Turing peut se résumer à la nécessité d'effectuer des choix sur les transitions,
mais qui n'ont aucun impact sur le résultat final.

En un sens, on supprime ainsi le non déterminisme, mais le non déterminisme est
souvent (plus qu'autre chose) un obstacle à l'implémentation pratique. Ainsi, on
est tenté d'exprimer la propriété pour un ARS que tout choix de transitions
mène au même résultat. Une première façon d'exprimer cette propriété est la
suivante.

\begin{definition}[Propriété de Church-Rosser]
    Soit $(A,\varTo)$ un ARS. On dit qu'il a la propriété de Church-Rosser 
    (abrégée CR) si, pour tous $a,b \in A$ tels que $a \varToRst b$, il existe
    $c \in A$ tel que $a \varToRt c$ et $b\varToRt c$.
\end{definition}

Pour tirer profit de cette propriété, introduisons aussi la notion de réécriture
correspondant à ce qu'on imagine être un résultat~: celle de forme normale.

\begin{definition}[Forme normale]
    Soit $(A,\varTo)$ un ARS. Un élément $a \in A$ est dit sous forme normale s'il
    n'existe pas de $b \in A$ tel que $a \varTo b$. On note $\NF(A,\varTo)$ (ou
    simplement $\NF$ quand le contexte est clair) l'ensemble des formes normales
    dans $(A,\varTo)$.
\end{definition}

L'intérêt de la propriété de Church-Rosser est qu'elle implique la propriété
suivante.

\begin{definition}[Propriété de la forme normale unique]
    Soit $(A,\varTo)$ un ARS. On dit qu'il a la propriété de la forme normale
    unique (\foreignexpr{Uniqueness of Normal Form} en anglais, abrégé en UNF)
    si pour toutes deux formes normales $a,b \in \NF(A,\varTo)$, si
    $a \varToRst b$ alors $a = b$.
\end{definition}

\begin{proposition}
    Pour tout ARS $(A,\varTo)$, si $(A,\varTo)$ a la propriété CR alors
    $(A,\varTo)$ a la propriété UNF.
\end{proposition}

\begin{proof}
    Soient $a,b\in \NF(A,\varTo)$ tels que $a \varToRst b$. Alors, par la
    propriété de Church-Rosser, on trouve $c$ tel que $a \varToRt c$ et
    $b\varToRt c$. Comme $a$ et $b$ sont des formes normales, on sait que
    $a \varToRt c$ (respectivement $b\varToRt c$) donne une réduction en $0$
    pas, c'est-à-dire que $a = c$ et $b = c$, d'où $a = b$.
\end{proof}

Ainsi, pour un système vérifiant CR, tout élément $a \in A$ peut se réduire en
au plus un élément.

\begin{notation}
    Soient $a,b \in \NF$. On note $a\varToMax b$ pour signifier que la réduction
    $a \varToRt b$ est maximale, et ainsi que $b \in \NF$.
\end{notation}

La détermination de formes normales est au c\oe ur de l'étude des systèmes de
calcul, et la théorie de la réécriture permet alors de donner des outils
efficaces pour prouver des propriétés sur ladite détermination dans le cas
général.

La suite de cette section est dédiée à la construction d'outils pour établir
la propriété CR. Pour ce faire, le premier pas est de donnée une formulation
alternative de la propriété CR, nommée confluence.

\begin{definition}[Confluence]
    Soit $(A,\varTo)$ un ARS. On dit qu'il est confluent si pour tous
    $a,b,c \in A$ tels que $a \varToRt b$ et $a \varToRt c$, il existe
    $d \in A$ tel que $b \varToRt d$ et $c\varToRt d$. De façon alternative,
    on peut exprimer ça par l'inclusion
    $\mathord{\varToTTr}\circ\mathord{\varToRt}\subseteq
    \mathord{\varToRt}\circ\mathord{\varToTTr}$.
\end{definition}

\begin{figure}[t!]
  \centering
  %\hrule

  %\vspace{0.2cm}
  \begin{tikzpicture}[node distance = 2cm, on grid, auto]
    \node (q_0) {$a$};
    \node (q_1) [below left = of q_0] {$b$};
    \node (q_2) [below right = of q_0] {$c$};
    \node (q_3) [below right = of q_1] {$d$};
    \node (r_1) [right = of q_2] {$b$};
    \node (r_0) [above right = of r_1] {$a$};
    \node (r_2) [below right = of r_0] {$c$};
    \node (r_3) [below right = of r_1] {$d$};
    \node (s_1) [right = of r_2] {$b$};
    \node (s_0) [above right = of s_1] {$a$};
    \node (s_2) [below right = of s_0] {$c$};
    \node (s_3) [below right = of s_1] {$d$};
    \node (a) at (0,-3.5) {Confluence};
    \node (b) at (5,-3.5) {Confluence locale};
    \node (c) at (10,-3.5) {Propriété du carreau};
    \draw[->,>=latex] (q_0) to node[very near end,below]{$\star$} (q_1);
    \draw[->,>=latex] (q_0) to node[very near end,below]{$\star$} (q_2);
    \draw[->,>=latex,dashed] (q_1) to node[very near end,below]{$\star$} (q_3);
    \draw[->,>=latex,dashed] (q_2) to node[very near end,below]{$\star$} (q_3);
    \draw[->,>=latex] (r_0) -- (r_1);
    \draw[->,>=latex] (r_0) -- (r_2);
    \draw[->,>=latex,dashed] (r_1) to node[very near end,below]{$\star$} (r_3);
    \draw[->,>=latex,dashed] (r_2) to node[very near end,below]{$\star$} (r_3);
    \draw[->,>=latex] (s_0) -- (s_1);
    \draw[->,>=latex] (s_0) -- (s_2);
    \draw[->,>=latex,dashed] (s_1) -- (s_3);
    \draw[->,>=latex,dashed] (s_2) -- (s_3);
  \end{tikzpicture}

  \vspace{0.2cm}
  \hrule

  \vspace{0.4cm}
  \caption{Confluence, confluence locale, propriété du carreau}
  \label{fig.confluence}
  %\hrule
\end{figure}

\begin{proposition}\label{prop.CR.confl}
    Tout ARS $(A,\varTo)$ a la propriété CR si et seulement s'il est confluent.
\end{proposition}

\begin{figure}[t]
  \centering
  %\hrule

  %\vspace{0.2cm}
  \begin{tikzpicture}[node distance = 2cm, on grid, auto]
    \node (a) {$a$};
    \node (bla) [right = of a] {};
    \node (c) [right = of bla] {$c$};
    \draw[<->,>=latex] (a) -- node[very near end, above right] {$\star$} (c);
    \node (b) [right = of c] {$b$};
    \draw[->,>=latex] (c) -- (b);
    \node (d0) [below = of bla] {$d_0$};
    \draw[->,>=latex,dashed] (a) -- node[very near end, right] {$\star$} (d0);
    \draw[->,>=latex,dashed] (c) -- node[very near end, below right] {$\star$} (d0);
    \node (d) [below right = of d0] {$d$};
    \draw[->,>=latex,dashed] (d0) -- node[very near end, right] {$\star$} (d);
    \draw[->,>=latex,dashed] (b) -- node[very near end, below right] {$\star$} (d);
  \end{tikzpicture}

  \vspace{0.2cm}
  \hrule
  
  \vspace{0.4cm}
  \caption{Preuve de la \cref{prop.CR.confl}}
  \label{fig.CR.confl}
  %\hrule

\end{figure}

\begin{proof}
    L'un des sens est évident~: si on a la propriété CR, alors (en reprenant les
    notations précédentes) $b \varToRst c$ donc on trouve un $d$ correspondant.

    On souhaite donc montrer la réciproque. Tout d'abord, on peut exprimer la
    relation $\varToRst$ comme étant $\rtClot{\varToS}$ (\latinexpr{c.f}
    l'\cref{exo.rstClot}). On suppose qu'on a la confluence et on souhaite
    montrer qu'on a la propriété CR. Soient $a,b \in A$ tels que $a\varToRst b$.
    On raisonne donc par induction sur le nombre de réécriture de $\varToS$
    pour prouver qu'il existe $c \in A$ tel que $a \varToRt c$ et $b\varToRt c$~:
    \begin{itemize}
    \item la propriété est évidente en $0$ pas puisqu'alors $a = b$~;
    \item on suppose donc que $a \varToRst c \varTo b$ et on souhaite trouver
      $d$ tel que $a \varToRt d$ et $b \varToRt d$. Par hypothèse d'induction,
      on trouve donc $d_0$ tel que $a \varToRt d_0$ et $c \varToRt d_0$. On
      sait donc que $c \varToRt d_0$ et $c\varToRt b$ (var $c \varTo b$), donc
      par confluence on en déduit qu'il existe $d$ tel que $d_0 \varToRt d$
      et $b\varToRt d$, d'où (par transitivité) que $a\varToRt d$.
    \item on suppose enfin que $a \varToRst c \varToTr b$ et on souhaite
      trouver $d$ tel que $a\varToRt d$ et $b \varToRt d$. Par hypothèse
      d'induction, on trouve $d$ tel que $a\varToRt d$ et $c\varToRt d$,
      mais alors par transitivité $b\varToRt d$, d'où le résultat.\qedhere
    \end{itemize}
\end{proof}

L'avantage de la confluence est qu'elle est une propriété plus naturelle à
vérifier~: on se donne deux réductions $\varToRt$ partant d'un même élément
$a$, et on trouve un point qui les relie. Cependant, cette propriété reste
globale, au sens où le nombre de cas à traiter est très grand. Le fait de
quantifier sur des triplets $b\varToTTr a \varToRt c$ donne des réductions
arbitrairement grandes. C'est pourquoi on introduit des variantes locales
de la confluence.

\begin{definition}[Confluence locale, propriété du carreau]
    Soit $(A,\varTo)$ un ARS. On dit qu'il a la confluence locale si pour tous
    $b \varToTr a \varTo c$, il existe $d \in A$ tel que
    $b \varToRt d \varToTTr c$.
    On dit qu'il a la propriété du carreau (ou la propriété du diamant, de
    l'anglais \foreignexpr{diamond property}) si pour tous
    $b \varToTr a \varTo c$ il existe $d \in A$ tel que
    $b \varTo d \varToTr c$.
\end{definition}

\begin{remark}
    La confluence pour $\varTo$ peut se voir exactement comme la propriété
    du carreau pour $\varToRt$.
\end{remark}

Un point essentiel dans ces propriétés est que la propriété du carreau
implique la confluence.

\begin{proposition}\label{prop.carreau.confluence}
    Soit $(A,\to)$ un ARS. Si $\varTo$ a la propriété du carreau, alors $\varTo$
    est confluente.
\end{proposition}

\begin{figure}[t!]
  \centering
  %\hrule

  %\vspace{0.2cm}
  \begin{tikzpicture}[node distance = 1.2cm, on grid, auto]
    \node (a) {$a$};
    \node (a_0) [below left = of a] {$a_0$};
    \node (a_1) [below right = of a] {$a_1$};
    \node (a') [below right = of a_0] {$a'$};
    \node (no) [below left = of a_0] {};
    \node (noo) [below right = of a_1] {};
    \node (b) [below left = of no] {$b$};
    \node (c) [below right = of noo] {$c$};
    \node (b_0) [below right = of b] {$b_0$};
    \node (c_0) [below left = of c] {$c_0$};
    \node (nooo) [below left = of c_0] {};
    \node (d) [below left = of nooo] {$d$};
    \draw[->,>=latex] (a) -- (a_0);
    \draw[->,>=latex] (a) -- (a_1);
    \draw[->,>=latex] (b) -- (b_0);
    \draw[->,>=latex] (c) -- (c_0);
    \draw[->,>=latex,dashed] (a_0) -- (a');
    \draw[->,>=latex,dashed] (a_1) -- (a');
    \draw[->,>=latex] (a_0) to node[very near end,below]{$p$} (b);
    \draw[->,>=latex] (a_1) to node[very near end,below]{$k$} (c);
    \draw[->,>=latex,dashed] (b_0) to node[very near end,below]{$k$} (d);
    \draw[->,>=latex,dashed] (c_0) to node[very near end,below]{$p$} (d);
    \draw[->,>=latex,dashed] (a') to node[very near end,below]{$p$} (b_0);
    \draw[->,>=latex,dashed] (a') to node[very near end,below]{$k$} (c_0);
  \end{tikzpicture}

  \vspace{0.2cm}
  \hrule
  
  \vspace{0.4cm}
  \caption{Preuve de la \cref{prop.carreau.confluence}}
  \label{fig.carreau.confluence}
  %\hrule
\end{figure}

\begin{proof}
  On prouve par récurrence sur $\max(p,k)$ que si $a\varTo^p b$ et
  $a\varTo^k c$ alors il existe $d$ tel que $b\varTo^{\leq k} d$ et
  $c\varTo^{\leq p} d$. Le cas où $p = 0$
  ou $k = 0$ est évident en prenant respectivement $d = c$ et $d = b$.

  On suppose donc la propriété vraie pour $n$, et on souhaite la
  prouver pour $n + 1$. Soit donc $a \varTo^{p+1} b$, $a\varTo^{k+1} c$.
  Par définition, on trouve $a_0,a_1$ tels que $a\varTo a_0 \varTo^p b$ et
  $a\varTo a_1 \varTo^k c$.

  Par propriété du diamant, on trouve donc $a'$ tel que
  $a_0 \varTo a'$ et $a_1 \varTo a'$. On peut alors appliquer l'hypothèse de
  récurrence à~:
  \begin{itemize}
  \item $a_0 \varTo^p b$, $a_0 \varTo a'$ pour trouver $b_0$ tel que
    $b \varTo b_0$ et $a' \varTo^{\leq p} b_0$.
  \item $a_1 \varTo^k c$, $a_1 \varTo a'$ pour trouver $c_0$ tel que
    $c \varTo c_0$ et $a' \varTo^{\leq k} c_0$.
  \end{itemize}
  mais alors, par hypothèse de récurrence sur
  $a' \varTo^{\leq p} b_0$, $a' \varTo^{\leq k} c_0$,
  on trouve $d$ tel que $b_0 \varTo^{\leq k} d$ et $c_0\varTo^{\leq p} d$,
  d'où $b \varTo^{\leq k+1} d$ et $c\varTo^{\leq p+1} d$. D'où le résultat par
  récurrence.
\end{proof}

On peut se demander alors si la confluence locale implique elle aussi la
confluence, mais cela n'est pas le cas en considérant le contre-exemple en
\cref{fig.ctrex.confluence}. On verra cependant un cas où la confluence
locale implique la confluence.

On termine cette section par une technique utile pour prouver la confluence, à
l'aulne des constructions précédentes.

\begin{lemma}\label{lem.carreau.confluence.clot}
    Soit $(A,\varTo)$ un ARS. Soit $\varRR$ une relation telle que~:
    \begin{itemize}
    \begin{minipage}{0.3\textwidth}
    \item $\mathord{\varRR}\subseteq\mathord{\varToRt}$~;
    \end{minipage}
    \begin{minipage}{0.3\textwidth}
    \item $\mathord{\varTo}\subseteq\mathord{\varRR}$~;
    \end{minipage}
    \begin{minipage}{0.3\textwidth}
    \item $\varRR$ a la propriété du carreau~;
    \end{minipage}
    \end{itemize}
    alors $\varTo$ est confluente.
\end{lemma}

\begin{proof}
    Tout d'abord, $\varRR$ contient $\varTo$, donc sa clôture réflexive et
    transitive est contenue dans celle de $\varTo$, et $\varTo$ contient
    $\varRR$~: puisque $\rtClot{\varRR}$ est la plus petite relation
    réflexive et transitive contenant $\varRR$, on en déduit l'inclusion
    inverse, d'où $\rtClot{\varRR} = \varToRt$. Comme $\varRR$ a la propriété
    du carreau, on en déduit que $\varRR$ est confluente, c'est-à-dire que
    $\rtClot{\varRR}$ a la propriété du carreau, ce que l'on peut par égalité
    appliquer à $\varToRt$ pour en déduire que $\varTo$ est confluente.
\end{proof}

\begin{figure}[t]
  \centering
  %\hrule

  %\vspace{0.2cm}
  \begin{tikzpicture}[node distance = 2cm, on grid, auto]
    \node (q_0) {$a$};
    \node (q_1) [right of = q_0] {$b$};
    \node (q_2) [right of = q_1] {$c$};
    \node (q_3) [right of = q_2] {$d$};
    \draw[->,>=latex] (q_1) -- (q_0);
    \draw[->,>=latex] (q_2) -- (q_3);
    \draw[->,>=latex] (q_1) to[bend left] (q_2);
    \draw[->,>=latex] (q_2) to[bend left] (q_1);
  \end{tikzpicture}

  \vspace{0.2cm}
  \hrule
  
  \vspace{0.4cm}
  \caption{Confluence locale sans confluence}
  \label{fig.ctrex.confluence}
  %\hrule
\end{figure}

On a donc une méthode efficace pour prouver la confluence, qui nous donne
une propriété forte impliquant que tout élément a une seule forme normale.
Intéressons-nous maintenant à la question de s'il existe toujours une telle
forme normale.

\subsection{Terminaison}

En effet, si la confluence est un outil important pour déduire que toute
forme normale représente de façon unique sa classe modulo $\varToRst$, il
est important de trouver aussi une telle forme normale.

On peut en fait trouver deux notions distinctes~: l'existence d'une forme
normale, d'une part, et la notion de normalisation forte.

\begin{definition}[Normalisation forte, faible]
    Soit $(A,\varTo)$ un ARS. On dit qu'un élément $a \in A$ est faiblement
    normalisant quand il existe $b \in \NF$ tel que $a\varToMax b$. On note
    $\WN(A,\varTo)$ l'ensemble des termes faiblement normalisant. On définit
    l'ensemble $\SN(A,\varTo)$ des termes fortement normalisant comme celui
    engendré par la règle
    \begin{prooftree}
        \AxiomC{$\forall b \in A, a \varTo b \implies b \in \SN$}
        \UnaryInfC{$a \in \SN$}
    \end{prooftree}
    On dit que l'ARS est faiblement normalisant quand $\WN = A$ et fortement
    normalisant quand $\SN = A$.
\end{definition}

La faible normalisation est probablement la définition la plus intuitive des
deux~: elle dit simplement que tout élément possède une forme normale (ainsi,
un système faiblement normalisant et confluent permet de parler de \og la\fg
forme normale d'un élément). Cependant, elle est au final moins manipulable,
car atteindre la forme normale d'un élément peut demander de choisir
précisément des réductions, ce qui peut devenir excessivement difficile (voire
incalculable). La forte normalisation, elle, peut se résumer à l'idée que
toute suite de réductions par $\varTo$ est finie, ce qui signifie qu'on
atteint toujours une forme normale en exécutant des réductions (évidemment,
le choix des réductions à effectuer peut influencer la taille du chemin vers
une forme normale).

\begin{exercise}
    Montrer que pour un ARS $(A,\varTo)$, l'ensemble $\SN$ est la plus grande
    partie de $A$ sur laquelle $\varTo$ est une relation bien fondée.
\end{exercise}

Sur la vision de $\SN$ comme une partie bien fondée, on peut deviner un outil
assez naturel pour montrer des résultats de normalisation (forte comme faible).

\begin{proposition}
    Soit $(A,\varTo)$ un ARS. Soit $(X,\varRR)$ un ensemble muni d'une relation
    bien fondée et $\makeFunName{f}{A}{X}$ une fonction. Si pour tous
    $a,b \in A$ tels que $a\varTo b$, $f(a)\varRR f(b)$, alors $\varTo$ est
    fortement normalisante. Si pour tout $a \in \NF\setminus A$ il existe
    $b \in A$ tel que $a\varTo b$ et $f(a)\varRR f(b)$ alors $\varTo$ est
    faiblement normalisante.
\end{proposition}

\begin{proof}
    Pour le cas de la forte normalisation, on raisonne par induction bien
    fondée sur $\varRR$. Soit donc $x \in X$ tel que pour tout $y\in X$ tel
    que $x \varRR y$ et tout $a \in A$, si $f(a) = y$ alors $a \in \SN$.
    Soit $a \in A$ tel que $f(a) = x$, alors pour tout $b \in A$ tel que
    $a \varTo b$, on sait que $f(a)\varRR f(b)$, donc par hypothèse
    d'induction $b \in \SN$, donc par définition $a \in \SN$. Donc pour
    tout $x \in X$, si $a \in A$ est tel que $f(a) = x$, alors $a \in \SN$,
    et puisque $f$ est une fonction totale, on en déduit que $\SN = A$.

    Pour le cas de la faible normalisation, on raisonne aussi par
    induction bien fondée sur $\varRR$. De même, soit $x \in X$ tel que pour
    tout $y \in X$ tel que $x \varRR y$ et tout $a \in A$, si $f(a) = y$
    alors il existe une forme normale à $a$. Soit alors $a$ tel que
    $f(a) = x$. Si $a$ est une forme normale, alors $a$ est faiblement
    normalisant. Sinon, alors on sait qu'il existe $b$ tel que
    $f(a)\varRR f(b)$, donc par hypothèse d'induction $b$ a une
    forma normale, et $a \varTo b \varToMax n$ avec $n$ une forme
    normale pour $b$, donc $a\varToMax n$, d'où le résultat par
    induction bien fondée.
\end{proof}

Ainsi, dans les cas les plus simples, il suffit de trouver une bonne mesure
décroissante lors d'une réécriture. Pour cela, on peut bien sûr associer des
entiers, mais on peut aussi utiliser la structure bien fondée des listes sur
un ensemble bien fondé, des paires avec l'ordre lexicographique\ldots

\begin{exercise}
    Soit un ARS $(A,\varTo)$, montrer que $\varTo$ est fortement normalisant
    si et seulement si $\varToT$ est fortement normalisant.
\end{exercise}

En particulier, la forte normalisation nous permet de raisonner par induction
sur $\varTo$. Pour prouver une propriété $P(x)$ sur un ARS $(A,\varTo)$
fortement normalisant, il nous suffit de montrer que pour tout $a \in A$, si
pour tout $b$ tel que $a \varTo b$, on a $P(b)$, alors on a $\varFF(a)$.
Cette méthode permet de prouver le lemme de Newman.

\begin{lemma}[Newman \cite{Newman42}]\label{lem.newman}
    Soit $(A,\varTo)$ un ARS. S'il est fortement normalisant et localement
    confluent, alors il est confluent.
\end{lemma}

\begin{proof}
    On procède, par induction sur $\varTo$, en considérant la propriété
    \[P(a) \defeq \forall b,c \in A, b \varToTTr a \varToRt c \implies
    \exists d \in A, b \varToRt d \varToTTr c\]
    On suppose donc donnés $a,b,c$ avec $b\varToTTr a \varToRt c$. Remarquons
    d'abord que si $b = a$ ou $c = a$, alors on peut prendre $d \defeq c$
    (respectivement $d \defeq b$) pour avoir le résultat. On suppose donc sans
    perte de généralité que $b\varToTTr b_0 \varToTr a \varTo c_0 \varToRt c$.
    On peut maintenant utiliser la confluence local en
    $b_0\varToTr a \varTo c_0$ pour trouver $d_0$ tel que
    $b_0 \varToRt d_0 \varToTTr c_0$, mais alors on peut appliquer l'hypothèse
    d'induction à $b\varToTTr b_0 \varToRt d_0$ pour trouver $d_1$ puis
    $d_1\varToTTr c_0\varToRt c$ pour en déduire le résultat.
\end{proof}

\begin{figure}[t]
    \centering
    \begin{tikzpicture}[node distance = 2cm, on grid, auto]
        \node (a) {$a$};
        \node (b0) [below left of = a] {$b_0$};
        \node (bla1) [below left of = b0] {};
        \node (b) [below left of = bla1] {$b$};
        \node (c0) [below right of = a] {$c_0$};
        \node (bla2) [below right of = c0] {};
        \node (c) [below right of = bla2] {$c$};
        \node (d0) [below right of = b0] {$d_0$};
        \node (d1) [below right of = b] {$d_1$};
        \node (bla3) [below right of = d1] {};
        \node (d) [below right of = bla3] {$d$};
        \draw[->,>=latex] (a) -- (b0);
        \draw[->,>=latex] (a) -- (c0);
        \draw[->,>=latex] (b0) -- node[very near end, below] {$\star$} (b);
        \draw[->,>=latex] (c0) -- node[very near end, right] {$\star$} (c);
        \draw[->,>=latex,dashed] (b0) -- node[very near end, right] {$\star$} (d0);
        \draw[->,>=latex,dashed] (c0) -- node[very near end, below] {$\star$} (d0);
        \draw[->,>=latex,dashed] (b) -- node[very near end, right] {$\star$} (d1);
        \draw[->,>=latex,dashed] (d0) -- node[very near end, below] {$\star$} (d1);
        \draw[->,>=latex,dashed] (d1) -- node[very near end, right] {$\star$} (d);
        \draw[->,>=latex,dashed] (c) -- node[very near end, below] {$\star$} (d);
    \end{tikzpicture}

    \vspace{0.2cm}
    \hrule

    \vspace{0.2cm}
    \caption{Preuve du \cref{lem.newman}}
\end{figure}

Un ARS qui est à la fois confluent et fortement normalisant est parfois appelé un
système terminant. Un point central dans un tel calcul est que, si on possède
un codage de Gödel $\godcod\tok$ de notre ARS $(A,\varTo)$ et que la relation de
réécriture est calculable sur les codes, alors le problème de décider si deux éléments
$a,b \in A$ vérifient $a\varToRst b$ est décidable.

Cette question, de décider si deux éléments sont égaux modulo réécriture,
a une importance historique toute particulière. Nous allons donc donner des
bases théoriques pour en comprendre l'importance.

\section{Problème du mot}

Pour énoncer le problème du mot, commençons par fixer un cadre plus restrictif à
notre étude. Les ARS constituent un cadre particulièrement général, si bien que
celui-ci empêche de dire grand chose. Aussi, nous préférons maintenant nous
restreindre à un cadre qui a déjà été introduit dans le \cref{chp.auto}~: celui
des monoïdes. Notre objectif va être de considérer des relations de réécriture
sur des monoïdes, en prenant comme idée centrale l'étude des monoïdes libres.

\subsection{Réécriture sur des monoïdes}

En effet, prenons par exemple le langage des formules $\Formula(\Sign)$. Comme
on considère un langage, on sait qu'on étudie une partie $\varLL\subseteq \toStB$
pour un certain alphabet $\ABB$ (qu'on pourrait ici expliciter en ajoutant les
symboles de fonctions et de relations aux connecteurs, quantificateurs, variables
\latinexpr{etc.}) qui se traduit donc par un certain monoïde libre. Considérons
par exemple la relation de réécriture suivante sur les termes arithmétiques~:
\begin{itemize}
    \begin{minipage}{0.45\textwidth}
    \item $t + 0 \reecr{} t$
    \end{minipage}
    \begin{minipage}{0.45\textwidth}
    \item $t + \Ss u \reecr{} \Ss(t + u)$
    \end{minipage}
\end{itemize}
cette relation induit en particulier une réécriture sur $\Formula(\Sign)$, mais
on peut tout de même l'étudier comme une relation sur $\toStB$. Cela permet en
particulier d'expliciter la notion de compatibilité, permettant de faire des
réécritures dans des sous-termes. Définissons tout cela formellement.

\begin{definition}[Relation compatible sur un monoïde]
    Soit $\varMMo$ un monoïde et $\varRR$ une relation sur $\varMMo$. On dit que
    $\varRR$ est compatible si
    \[\forall m,m' \in \varMMo, m \varRR m' \implies
    \forall m_{\rml}, m_{\rmr} \in \varMMo,
    m_{\rml}\opM m\opM m_{\rmr} \varRR m_{\rml} \opM m' \opM m_{\rmr}\]
    On dit que $\varRR$ est une congruence si, de plus, $\varRR$ est une
    relation d'équivalence.
\end{definition}

Ainsi, les relations qui vont nous intéresser sont les relations compatibles,
mais la donnée d'une relation compatible en elle-même signifie la donnée d'une
infinité d'instances (pour peu que le monoïde est infini). On souhaite donc
construire un cadre dans lequel on peut trouver des données finies résumant
notre relation. Pour cela, on introduit la clôture compatible d'une relation.

\begin{definition}[Clôture compatible d'une relation]
    Soit $\varRR$ une relation sur un monoïde $\varMMo$. On appelle clôture
    compatible la plus petite relation $\toCompt{\varRR}$ contenant $\varRR$
    et compatible. Dans le cas d'un ARS $(\varMMo,\varTo)$, on notera
    $\varToC$ la clôture compatible de $\varTo$, ainsi que
    $\varToCRst$ la plus petite congruence contenant $\varTo$.
\end{definition}

Reprenons maintenant notre exemple précédent de $\reecr{}$ sur les formules.
La donnée de la relation est uniquement celle de couples de la forme
\[\compr{(``t+0",``t")}{t \in \Term(\SignP{\Arith})}\cup
\compr{(``t+\Ss u",``\Ss(t+u)"}{t,u \in \Term(\SignP{\Arith})}\]
mais, en considérant la relation compatible qu'elle engendre, on peut
maintenant écrire par exemple
\[\Ss(\Ss 0 + \Ss 0) = \Ss\Ss\Ss 0 \rtClot{\reecr{}}
\Ss\Ss\Ss 0 = \Ss\Ss\Ss 0\]
Nous repoussons à la fin du chapitre la discussion de systèmes de preuves
dans lequel on autorise ce genre de réductions, mais il est clair ici que
la seconde proposition est largement plus simple à prouver (elle se prouve
par simple réflexivité de l'égalité).

En nous restreignant au cas des monoïdes, ces notions nous permettent de
construire ce qu'on appelle une présentation de monoïde.

\begin{definition}[Présentation de monoïde]
    On appelle présentation de monoïde une paire $(\ABA,\varTo)$ où
    $\ABA$ est un ensemble quelconque et
    $\mathord{\varTo}\subseteq\left(\toStA\right)^2$. Le monoïde
    engendré par une présentation $(\ABA,\varTo)$ est le monoïde
    quotient $\quot{\toStA}{\varToCRst}$. Un monoïde $\varMMo$
    admet une présentation $(\ABA,\varTo)$ dont le monoïde engendré
    est isomorphe à $\varMMo$. On dit que $\varMMo$ est finiment
    présenté s'il existe une présentation dont $\ABA$ et
    $\varTo$ sont des ensembles finis, et finiment engendré s'il
    existe une présentation dont $\ABA$ est un ensemble fini.
\end{definition}

On peut maintenant énoncer le problème du mot, qui s'applique à des
présentations de monoïdes.

\begin{definition}[Problème du mot]
    Soit $(\ABA,\varTo)$ une présentation de monoïde et un codage de Gödel
    $\makeFunName{\godcod\tok}{\ABA}{\bN}$ (qu'on étend sans souci à
    $\toStA$ par la bijection $\BijCs$). Le problème du mot pour cette
    présentation est la question de savoir si l'ensemble
    \[\compr{(\godcod a,\godcod b)}{a,b \in \toStA\land a\varToCRst b}\]
    est décidable.
\end{definition}

On peut aussi restreindre ce problème à une sous-classe des monoïdes, qui
est celle des groupes. On se réfère par exemple à \cite{algGrandCombat}
pour un cours plus détaillé abordant cette notion.

Le fait important étant que, même dans le cas où l'on considère simplement
les groupes, le problème du mot est indécidable. Ce théorème est donné
avant tout pour la culture, nous n'en donnons pas la preuve ici.

\begin{theorem}[Boone-Novikov \cite{novikov55, boone58}]
    Il existe un groupe finiment présenté dont le problème du mot est
    indécidable.
\end{theorem}

\subsection{Complétion de Knuth-Bendix}

Si le problème du mot en toute généralité est indécidable, même pour des
classes restreintes, il reste décidable dans de nombreux cas, et il est
alors intéressant d'étudier la relation de réécriture. Nous avons déjà vu
que, grâce au \cref{lem.newman}, la forte normalisation permet de se
contenter de la confluence locale pour prouver la confluence. Supposons
donc maintenant qu'on dispose d'un ARS $(\varMMo,\varTo)$ sur un monoïde,
fortement normalisant, dont on veut établir la confluence locale. L'étude
de cette confluence locale peut en fait se réduire à celle des paires
critiques, qui sont des conflits de règles. Prenons un exemple~:
soit le système de réécriture sur $\{\varla,\varlb\}$ donné par
\[\varla\varla \varTo \Ewd \qquad
\varlb\varlb \varTo \Ewd\qquad
\varlb\varla\varlb\varTo\varla\varlb\varla\]
Si on considère un mot $\varwd$ qui se réécrit en deux mots
$\varwv_1,\varwv_2$, on peut avoir~:
\begin{itemize}
\item le cas où $\varwd = \varwd_1\varwd_2$, $\varwv_1=\varwv_1'\varwd_2$
  et $\varwd_1\varTo\varwv_1'$, $\varwv_2 = \varwd_1\varwv_2'$ et
  $\varwd_2\varTo\varwv_2'$~;
\item le cas où $\varwv_1 = \varwv_2$~;
\item le cas où $\varwd$ contient $\varla\varla\varla$ et la réécriture
  réduit les deux premiers $\varla$ pour $\varwd_1$, et les deux derniers
  pour $\varwd_2$~;
\item la même chose avec $\varlb\varlb\varlb$~;
\item le cas où $\varwd$ contient $\varlb\varla\varlb\varla\varlb$
  et la réécriture réduit le $\varlb\varla\varlb$ de gauche pour donner
  $\varla\varlb\varla\varla\varlb$ d'un côté, et réduit le
  $\varlb\varla\varlb$ de droite pour donner
  $\varlb\varla\varla\varlb\varla$ de l'autre~;
\item le cas où $\varwd$ contient $\varlb\varlb\varla\varlb$~;
\item le cas où $\varwd$ contient $\varlb\varla\varlb\varlb$.
\end{itemize}
Comment prouver la confluence locale de ce système~? D'abord, les 2 premiers
cas peuvent se \og rejoindre\fg directement, le premier en effectuant
l'autre réduction puisque les deux peuvent se faire en quelque sorte en
parallèle, et le second en remarquant que $\varToRt$ contient la réflexivité,
donc le résultat est déjà vérifié.

Il ne reste donc qu'à vérifier les 5 autres cas, qu'on peut tous lister en
considérant directement la forme des mots mis en cause dans la définition de
$\varTo$. En effet, on a simplement cherché toutes les paires de mots
$(\varwd_1,\varwd_2)$ telles qu'un suffixe de $\varwd_1$ est aussi un préfixe
de $\varwd_2$, et $\varwd_1$ et $\varwd_2$ sont des membres droits de $\varTo$.
Ainsi, en posant $\varwd_1 = \varwd_1'\varww$ et $\varwd_2 = \varww\varwd_2'$,
le mot $\varwd_1'\varww\varwd_2'$ est le mot associé à la paire critique
$(\varwd_1,\varwd_2)$. C'est cette étude qui va permettre de traiter de façon
automatique les systèmes de réécriture.

\begin{definition}[Paire critique]
    Soit $(\ABA,\varTo)$ une présentation de monoïde. On appelle paire
    critique sur cette présentation une paire $(\varwd_1,\varwd_2)$
    telle que
    \[\exists \varwd_1',\varwd_2',\varwv_1,\varwv_2,\varww\in\toStA,\begin{cases}
        \varwd_1=\varwd_1'\varww\\
        \varwd_2=\varww\varwd_2'\\
        \varwd_1\varTo \varwv_1\\
        \varwd_2\varTo \varwv_2
    \end{cases}\]
    On dit que le mot associé à cette paire est, dans la décomposition
    donnée au-dessus, le mot $\varwd_1'\varww\varwd_2'$. On dit que cette
    paire peut être rejointe (ou résolue) s'il existe $\varwv$ tel que
    $\varwv_1\varToCRt \varwv\varToCTTr \varwv_2$.
\end{definition}

\begin{remark}
    La relation $\varToC$ est localement confluente si et seulement si toute
    paire critique peut être rejointe.
\end{remark}

Dans notre exemple, on peut remarquer que la réduction est fortement
normalisante~: en associant à chaque mot $\varwd$ la paire ordonnée
lexicographiquement $(\wul_{\varlb},\wul_{\varla})$, chaque règle fait
décroître strictement la paire.

Cependant, la confluence locale n'est ici pas assurée. Les paires
critiques $(\varla\varla,\varla\varla)$ et $(\varlb\varlb,\varlb\varlb)$
se résolvent de façon évidente, mais la troisième paire critique,
$\varlb\varla\varlb\varla\varlb$, donne deux mots incompatibles avec
$\varla\varlb\varla\varla\varlb$ (qui peut se réduire uniquement en
$\varla\varlb\varlb$) et $\varlb\varla\varla\varlb\varla$ (qui peut se
réduire uniquement en $\varlb\varlb\varla$). Les deux dernières paires
aussi sont impossible à résoudre en l'état.

On introduit donc un algorithme permettant, pour une présentation de monoïde,
de réécrire la relation $\varTo$ de sorte à obtenir une présentation
confluente. Pour cela A FAIRE